\documentclass[11pt,a4paper]{article}
\usepackage[margin=18mm]{geometry}
\usepackage{graphicx}
\usepackage{hyperref}
\usepackage{caption}
\hypersetup{colorlinks=true,urlcolor=blue,linkcolor=blue}
\title{Reproducing GitHub Insights for Diffusers}
\author{26AI60R04\\CD61203: Essential Tools for Scientific Computing}
\date{13 November 2026}
\newcommand{\comparison}[1]{%
  \begin{minipage}[t]{0.49\linewidth}\centering
    \textbf{Matplotlib reproduction}\par\smallskip
    \includegraphics[width=\linewidth]{#1/generated.png}
  \end{minipage}\hfill
  \begin{minipage}[t]{0.49\linewidth}\centering
    \textbf{GitHub reference}\par\smallskip
    \includegraphics[width=\linewidth]{#1/reference.png}
  \end{minipage}}
\begin{document}
\maketitle
\noindent The four figures reproduce charts from the public repository
\href{https://github.com/huggingface/diffusers}{huggingface/diffusers}.
References were captured from GitHub on 13 November 2026, in the light theme.
The data, plotted values, source URLs, capture times, and browser settings are
retained in the repository. Each reproduction is drawn with Python/Matplotlib
from CSV, independently of the reference screenshot.

\begin{figure}[htbp]
\centering
\comparison{fig1-contributors}
\caption{Contributors: the first ten ranked contributor cards for the default
``Last 3 months'' selection, with 14 weekly bars per author, from 9 August to
8 November 2026. Blue bars show commits to \texttt{main}, excluding merge
commits. The right-hand axes use the common scale of the reference. Names,
rank, commit totals, additions, and deletions follow GitHub's displayed order,
including ties. The source contains additional contributors; the matching
reference region and reproduction both show the first ten.}
\label{fig:contributors}
\end{figure}
\noindent\small Source:
\url{https://github.com/huggingface/diffusers/graphs/contributors}.

\clearpage
\section*{Weekly commits}
\begin{figure}[htbp]
\centering
\comparison{fig2-commits}
\caption{Commits over the last year: 52 weekly totals from 16 November 2025
through 8 November 2026. The green column chart preserves the month labels,
the ``Commits'' axis label, and the 0--60 vertical scale. The largest weekly
total is 55, in the week beginning 4 January 2026.}
\label{fig:commits}
\end{figure}
\noindent Source:
\url{https://github.com/huggingface/diffusers/graphs/commit-activity}.

\medskip\noindent\textbf{Data and method.}
The GitHub chart endpoint supplies a total and seven daily counts for every
week. The CSV stores all eight counts and the UTC week-start date. The figure
uses the supplied weekly total without smoothing or interpolation. As an
independent consistency check, every total equals the sum of its daily counts.

\medskip\noindent\textbf{Comparison.}
All 52 values agree exactly with the accessible point labels captured from
GitHub's rendered chart. Font rendering, rounded corners, and interactive
controls are not required to be pixel-identical.

\clearpage
\section*{Code frequency}
\begin{figure}[htbp]
\centering
\comparison{fig3-code-frequency}
\caption{Code frequency across repository history. Green columns are additions;
red columns are signed deletions below zero. The raw endpoint provides 233
weekly observations, from 29 May 2022 to 8 November 2026. At the captured
full-history view GitHub groups those observations into 55 calendar-month
sums; the reproduction uses those same sums to match the visible reference.
The title and subtitle retain GitHub's wording, including ``per week''.}
\label{fig:code-frequency}
\end{figure}
\noindent Source:
\url{https://github.com/huggingface/diffusers/graphs/code-frequency}.

\medskip\noindent\textbf{Aggregation.}
Each weekly observation is assigned to the calendar month containing its UTC
week-start date. Additions and signed deletions are summed separately. The
weekly source remains in \texttt{raw\_data.csv}; the monthly values are also
exported as \texttt{displayed\_data.csv}. Plotting code recomputes the sums from
the weekly source. All 110 displayed series values agree exactly with the
captured GitHub point labels. The largest monthly additions count is 76,969;
the most negative monthly deletions count is $-46,931$.

\clearpage
\section*{Top committers}
\begin{figure}[htbp]
\centering
\comparison{fig4-top-committers}
\caption{The actual ``Top committers'' chart in GitHub Insights/Pulse for the
displayed period 5--12 November 2026. All 15 authors appear in the original
order. The green columns, 0--40 ``Commits'' scale, and avatar category labels
match the reference. The leading counts are 24 for sayakpaul, 19 for DN6,
18 for yiyixuxu, and 16 for claude.}
\label{fig:top-committers}
\end{figure}
\noindent Source: \url{https://github.com/huggingface/diffusers/pulse}.

\medskip\noindent\textbf{Scope.}
This chart is located inside Pulse rather than in a separate Top committers
menu entry. Its period differs from the Contributors page. Values are taken
directly from the rendered chart's accessible point labels; author names are
retained in the CSV even though GitHub displays avatars on the category axis.
The reference capture waits until all 15 avatar labels have loaded successfully.
Pulse's prose summary reports separate counts excluding merges; those summary
counts are not substituted for the plotted author counts.

\medskip\noindent\textbf{Reproducibility and validation.}
Run \texttt{python reproduce.py} to regenerate all figures offline, and
\texttt{python tools/validate.py} to compare all 317 plotted values against
the saved GitHub reference labels. Run \texttt{pdflatex main.tex} twice, or
\texttt{tectonic main.tex}, to build this document. The comparisons in
Figures~\ref{fig:contributors}--\ref{fig:top-committers} use authentic browser
screenshots of the corresponding chart regions. No synthetic reference images
or invented data are used.
\end{document}
